%!TEX root = ./main.tex

% This deck uses \text{...} inside math and \gtrsim; the shared variables.tex does
% not load the AMS packages, so load them here.
\usepackage{amsmath,amssymb}

\usetikzlibrary{positioning,fit,backgrounds}
\usetikzlibrary{shapes.geometric,calc}

\definecolor{ranblue}{RGB}{46,91,186}
\definecolor{coreorange}{RGB}{222,122,36}
\definecolor{ranpurple}{RGB}{109,53,151}
\definecolor{softgreen}{RGB}{73,139,74}
\definecolor{softgray}{RGB}{92,98,105}

\newcommand{\source}[2]{%
  \par\vspace{0.25em}%
  {\tiny\textcolor{gray}{\href{#1}{\underline{#2}}}\par}%
}

% Inline expansion of an acronym at first use: SMS \expand{Short Message Service}
\newcommand{\expand}[1]{\textcolor{softgray}{(#1)}}

% Small gray glossary line, for slides with several acronyms at once.
\newcommand{\glossline}[1]{%
  \par\vspace{0.4em}%
  {\scriptsize\color{softgray}\raggedright#1\par}%
}

\newcommand{\plainterm}[2]{%
  \textbf{#1}\hspace{0.25em}{\small\textcolor{softgray}{#2}}%
}

\newcommand{\ask}[1]{%
  \begin{tcolorbox}[colback=blue!4,colframe=ranblue,boxrule=0.7pt,arc=2pt,
    left=5pt,right=5pt,top=4pt,bottom=4pt]
    \centering\large #1
  \end{tcolorbox}%
}

% Font is deliberately \normalsize, not \large: at \large a one-sentence takeaway
% wraps to two lines and the box grows into whatever is below it.
\newcommand{\takeaway}[1]{%
  \begin{tcolorbox}[colback=orange!7,colframe=coreorange,boxrule=0.7pt,arc=2pt,
    left=5pt,right=5pt,top=4pt,bottom=4pt]
    \centering\normalsize #1
  \end{tcolorbox}%
}

% Beamer overlays inside TikZ: \visible<...> breaks path parsing inside node text,
% so reveal nodes with the standard opacity idiom instead: \node[visible on=<2->] ...
\tikzset{
  invisible/.style={opacity=0,text opacity=0},
  visible on/.style={alt=#1{}{invisible}},
  alt/.code args={<#1>#2#3}{\alt<#1>{\pgfkeysalso{#2}}{\pgfkeysalso{#3}}},
}

\tikzset{
  netnode/.style={draw=ranblue,fill=blue!5,rounded corners=3pt,
    minimum width=1.9cm,minimum height=0.82cm,align=center,font=\bfseries},
  corenode/.style={draw=coreorange,fill=orange!7,rounded corners=3pt,
    minimum width=1.9cm,minimum height=0.82cm,align=center,font=\bfseries},
  oranode/.style={draw=ranpurple,fill=purple!6,rounded corners=3pt,
    minimum width=1.9cm,minimum height=0.82cm,align=center,font=\bfseries},
  flow/.style={-{Latex[length=2.4mm]},thick,draw=softgray}
}

% Full-width question/answer block: spans the whole frame, no box.
% Use below \end{columns} when the question is meant for the whole slide.
\newcommand{\qa}[2]{%
  \par\vspace{0.55em}%
  {\large\raggedright
    \noindent\hangindent=1.25em\hangafter=1
    \textcolor{ranblue}{\textbf{Q:}}\hspace{0.35em}\textbf{#1}\par
    \vspace{0.35em}
    \noindent\hangindent=1.25em\hangafter=1
    \textcolor{coreorange}{\textbf{A:}}\hspace{0.35em}#2\par}%
}

% ---------------------------------------------------------------------------
% Class 12 (Wi-Fi Sensing) additions.
% COLOR CONVENTION for this deck -- every figure obeys it:
%   ranblue    = the ANTENNA axis (space)      -> reading it gives angle
%   coreorange = the SUBCARRIER axis (frequency) -> reading it gives delay/distance
%   softgreen  = the direct path (the answer we want)
%   ranpurple  = reflected paths (multipath ghosts)
%   softgray   = hardware dirt: clocks, offsets, detection delay
%   talkred    = bad news (merged taps, broken assumptions)
% ---------------------------------------------------------------------------

\definecolor{talkred}{RGB}{163,26,52}

% One antenna element: a small triangle on a mast, pointing up. #1 = (x,y), #2 = color
\newcommand{\anttri}[2]{%
  \draw[#2,thick] ($#1+(-0.11,0.14)$) -- ($#1+(0,-0.06)$) -- ($#1+(0.11,0.14)$)
    -- cycle;
  \draw[#2,thick] ($#1+(0,-0.06)$) -- ($#1+(0,-0.22)$);%
}

% The CSI matrix: 3 antennas x 8 subcarriers, one phasor arrow per cell. The phase
% step per column (26 deg) and per row (90 deg = the running example's quarter turn)
% are EXAGGERATED on the frequency axis so students can see the twist; every frame
% that shows it says so in a gray caption. Draw inside a tikzpicture at the origin.
\newcommand{\csigrid}{%
  \foreach \m in {0,1,2}{
    \foreach \k in {0,...,7}{
      \draw[softgray!50,fill=gray!6] (\k*0.72,\m*0.72)
        rectangle ++(0.6,0.6);
      \pgfmathsetmacro{\ang}{100-26*\k-90*\m}
      \draw[-{Latex[length=1.3mm]},thick,black!70]
        (\k*0.72+0.3,\m*0.72+0.3) -- ++(\ang:0.21);
    }}
  % axis labels
  \draw[-{Latex[length=2mm]},thick,coreorange] (0,-0.42) -- (5.9,-0.42);
  \node[font=\scriptsize\bfseries,text=coreorange] at (2.95,-0.75)
    {subcarriers $f_1 \ldots f_K$ (frequency)};
  \draw[-{Latex[length=2mm]},thick,ranblue] (-0.42,0) -- (-0.42,2.2);
  \node[font=\scriptsize\bfseries,text=ranblue,rotate=90] at (-0.78,1.1)
    {antennas $1 \ldots M$};%
}
